%%% Vyplňte prosím základní údaje o závěrečné práci (odstraňte {...}).
%%% Automaticky se pak vloží na všechna místa, kde jsou potřeba.

% Druh práce:
%	"bc" pro bakalářskou
%	"mgr" pro diplomovou
%	"phd" pro disertační
%	"rig" pro rigorozní
\def\ThesisType{mgr}

% Název práce v jazyce práce (přesně podle zadání)
\def\ThesisTitle{{Rozbíjení symetrií pro řešení splnitelnosti}}

% Název práce v angličtině
\def\ThesisTitleEN{{Symmetry Breaking Constraints for SAT Solving}}

% Jméno autora (vy)
\def\ThesisAuthor{{Martin Struk}}

% Rok odevzdání
\def\YearSubmitted{{2026}}

% Název katedry nebo ústavu, kde byla práce oficiálně zadána
% (dle Organizační struktury MFF UK:
% https://www.mff.cuni.cz/cs/fakulta/organizacni-struktura,
% případně plný název pracoviště mimo MFF)
\def\Department{{Katedra algebry}}
\def\DepartmentEN{{Department of algebra}}

% Jedná se o katedru (department) nebo o ústav (institute)?
\def\DeptType{{Katedra}}
\def\DeptTypeEN{{Department}}

% Vedoucí práce: Jméno a příjmení s~tituly
\def\Supervisor{{Mgr. Mikoláš Janota, Ph.D.}}

% Pracoviště vedoucího (opět dle Organizační struktury MFF)
\def\SupervisorsDepartment{{Katedra algebry}}
\def\SupervisorsDepartmentEN{{Department of algebra}}

% Studijní program (kromě rigorozních prací)
\def\StudyProgramme{{Matematika pro informační technologie}}

% Nepovinné poděkování (vedoucímu práce, konzultantovi, tomu, kdo
% vám nosil pizzu a vařil čaj apod.)
\def\Dedication{%
{Chtěl bych poděkovat vedoucímu práce Mgr. Mikoláši Janotovi, Ph.D.
za všechny rady, připomínky, vstřícnost, ochotu a motivaci během vedení této
práce.}
}

% Abstrakt (doporučený rozsah cca 80-200 slov; nejedná se o zadání práce)
\def\Abstract{%
{Symetrie v problémech booleovské splnitelnosti (SAT) často způsobují nadbytečné prohledávání ekvivalentních oblastí prohledávacího prostoru, což výrazně snižuje výkon řešičů. Klauzule pro rozbíjení symetrií tento problém řeší eliminací symetrických řešení při zachování splnitelnosti. V této práci se zabýváme generováním těchto klauzulí ze známých symetrií pro danou SAT formuli. Představujeme nástroj pro konstrukci klauzulí pro rozbíjení symetrie s využitím vzorů, což jsou kompaktní reprezentace zavedené v předchozí práci. Tento přístup generuje klauzule pomocí vybraných množin permutací a zároveň umožňuje řídit hloubku generování klauzulí, čímž umožňuje volit kompromis mezi silou rozbíjení symetrií a velikostí výsledné formule. Hlavní důraz je kladen na experimentální vyhodnocení toho, jak různé množiny permutací ovlivňují účinnost rozbíjení symetrií. Experimenty jsou prováděny na problémech, které lze přirozeně reprezentovat jako binární matice, avšak je možné uvažovat i složitější problémy. Hodnotíme jak výkon řešení, tak redukci počtu symetrických řešení pomocí moderních SAT řešičů.}
}

% Anglická verze abstraktu
\def\AbstractEN{%
{Symmetries in Boolean satisfiability (SAT) problems often cause redundant exploration of equivalent regions of the search space, significantly degrading solver performance. Symmetry breaking constraints (SBCs) address this issue by eliminating symmetric solutions while preserving satisfiability. In this work, we study the generation of SBCs from known symmetries for a given SAT formula. We present a tool for constructing SBCs using patterns, a compact representation introduced in previous work. The proposed approach generates clauses using selected sets of permutations while allowing control over the depth of generating constraints, providing a trade-off between the strength of symmetry breaking and the size of the resulting formula. Our main focus is an experimental evaluation of how different sets of permutations influence the effectiveness of symmetry breaking. The experiments are conducted on problems naturally representable as binary matrices, but more complex problems may be considered. We evaluate both solving performance and the reduction in the number of symmetric solutions using state-of-the-art SAT solvers.}
}

% 3 až 5 klíčových slov (doporučeno) oddělených \sep
% Hodí se pro nalezení práce podle tématu.
\def\ThesisKeywords{%
{booleovská splnitelnost \sep rozbíjení symetrií \sep generování klauzulí \sep permutační grupy}
}

\def\ThesisKeywordsEN{
{boolean satisfiability \sep symmetry breaking \sep constraint generation \sep permutation groups}
}

% Pokud některá z položek metadat obsahuje TeXové řídící sekvence, je potřeba
% dodat i verzi v obyčejném textu, která se objeví v metadatech formátu XMP
% zabudovaných do výstupního souboru PDF. Pokud si nejste jistí, prohlédněte si
% vygenerovaný soubor thesis.xmpdata.
\def\ThesisAuthorXMP{\ThesisAuthor}
\def\ThesisTitleXMP{\ThesisTitle}
\def\ThesisKeywordsXMP{\ThesisKeywords}
\def\AbstractXMP{\Abstract}

% Máte-li dlouhý abstrakt a nechceme se mu vejít na informační stranu,
% můžete použít toto nastavení ke zmenšení písma informační strany.
% (Uvažte nicméně zkrácení abstraktu, to je často lepší.)
%\def\InfoPageFont{}
\def\InfoPageFont{\small}  % odkomentujte pro zmenšení písma
